\documentclass[11pt]{article}
\usepackage{amsmath,amsthm,amssymb}
\usepackage[margin=1.2in]{geometry}
\usepackage{booktabs}

\newtheorem{theorem}{Theorem}[section]
\newtheorem{lemma}[theorem]{Lemma}
\newtheorem{proposition}[theorem]{Proposition}
\newtheorem{corollary}[theorem]{Corollary}
\newtheorem{conjecture}[theorem]{Conjecture}
\theoremstyle{definition}
\newtheorem{definition}[theorem]{Definition}
\theoremstyle{remark}
\newtheorem{remark}[theorem]{Remark}

\newcommand{\ZZ}{\mathbb{Z}}
\newcommand{\QQ}{\mathbb{Q}}
\newcommand{\CC}{\mathbb{C}}
\newcommand{\Sym}{\mathfrak{S}}
\newcommand{\im}{\operatorname{im}}
\newcommand{\rk}{\operatorname{rank}}

\title{Block-Multiplicative Structure of the Staircase Determinant\\
\large\normalfont\itshape (Partial Result --- Computational Evidence)}
\author{Clio}
\date{April 17, 2026}

\begin{document}
\maketitle

\begin{abstract}
Let $\Pi = \Pi_{n-1}$ be the staircase product in the group algebra
$\QQ[\Sym_n]$, let $D \subset \Sym_n$ be the set of descent-paired
permutations, and let $M_R[D,D]$ be the right-multiplication matrix
with entries $f_{w^{-1}v}$ (coefficient of $w^{-1}v$ in $\Pi$).
We document computational evidence for a \emph{block-multiplicative}
structure of $\det(M_R[D,D])$: as the staircase is built one block at
a time, the determinant multiplies by a strictly positive integer
$R_k$ at each step.  This would give $\det(M_R[D,D]) > 0$ by
induction, completing the $H$-invariant theorem.

The block structure is supported by exact computation for $n \le 6$.
No proof is given.  We record the precise factored values of each
ratio $R_k$, describe the tensor product phenomenon explaining
$R_2 = R_3 = 24^{|D_n|/|D_4|}$, and document the failure of the
most natural approach (total non-negativity).
\end{abstract}


%% ====================================================================
\section{Setup}
%% ====================================================================

\subsection{The symmetric group and its generators}

Let $\Sym_n$ denote the symmetric group on $n$ letters, with simple
transpositions $s_1, s_2, \ldots, s_{n-1}$ satisfying the Coxeter
relations:
\begin{align*}
  s_i^2 &= e,\\
  s_i s_j &= s_j s_i \quad (|i-j| \ge 2),\\
  s_i s_{i+1} s_i &= s_{i+1} s_i s_{i+1}.
\end{align*}
The longest element of $\Sym_n$ is denoted $w_0$; it has length
$\ell(w_0) = \binom{n}{2}$.

\subsection{The staircase word and staircase product}

\begin{definition}[Staircase word]
The \emph{staircase word} for $w_0 \in \Sym_n$ is the sequence of
simple transpositions
\[
  \mathbf{i} \;=\;
  \underbrace{s_1}_{k=1},\;
  \underbrace{s_2,\, s_1}_{k=2},\;
  \underbrace{s_3,\, s_2,\, s_1}_{k=3},\;
  \ldots,\;
  \underbrace{s_{n-1},\, s_{n-2},\, \ldots,\, s_1}_{k=n-1}.
\]
We call the portion for fixed $k$ \emph{block $k$}.  Block $k$ consists
of the $k$ generators $(s_k, s_{k-1}, \ldots, s_1)$ in that order.
\end{definition}

\begin{definition}[Staircase product]
The \emph{staircase product} is
\[
  \Pi \;=\; \prod_{j \in \mathbf{i}} \frac{1 + s_j}{2}
  \;\in\; \QQ[\Sym_n],
\]
where the product is taken left-to-right following the staircase word.
Each factor $P_j = (1+s_j)/2$ is a projection.  The \emph{partial
staircase product after $k$ blocks} is
\[
  \Pi_k \;=\; \prod_{j=1}^{k} \prod_{i=j}^{1} \frac{1+s_i}{2},
\]
so that $\Pi_{n-1} = \Pi$.
\end{definition}

The image $\Pi_k(\QQ[\Sym_n])$ stabilizes to the subgroup algebra
$\QQ[\Sym_{k+1}] \subset \QQ[\Sym_n]$ in a suitable sense: $\Pi_k$
is the staircase element for $\Sym_{k+1} \hookrightarrow \Sym_n$
acting from the left on the first $k+1$ positions.

\subsection{Descent-paired permutations}

\begin{definition}[Descent-paired permutations]
A permutation $w \in \Sym_n$ is \emph{descent-paired} if
$w(2i-1) > w(2i)$ for all $i = 1, 2, \ldots, \lfloor n/2 \rfloor$.
Let $D = D_n$ denote the set of all descent-paired permutations.
\end{definition}

The cardinality $|D_n| = n!/\,2^{\lfloor n/2 \rfloor}$.  The first
few values are:
\[
  |D_2| = 1,\quad |D_3| = 3,\quad |D_4| = 6,\quad
  |D_5| = 30,\quad |D_6| = 90.
\]
The set $D_n$ is the right coset representative system for the
subgroup $H = \langle s_1, s_3, s_5, \ldots \rangle \le \Sym_n$,
and $\Sym_n = D_n \cdot H$ (disjoint union of left cosets).

\subsection{The right-multiplication matrix}

Write $\Pi = \sum_{g \in \Sym_n} f_g\, g$ where $f_g \ge 0$ for all~$g$
(every permutation appears with non-negative coefficient in the staircase
product).  The coefficients satisfy the palindromic symmetry
$f_g = f_{g^{-1}}$ for all $g$.

\begin{definition}[Right-multiplication matrix]
The \emph{right-multiplication matrix} $M_R[D,D]$ is the
$|D| \times |D|$ matrix with rows and columns indexed by $D$, and
\[
  M_R[v, w] \;=\; f_{w^{-1} v}, \qquad v, w \in D.
\]
Equivalently, $M_R[v,w]$ is the coefficient of~$v$ in $w \cdot \Pi$,
restricted to $D$.
\end{definition}

By the palindromic symmetry $f_g = f_{g^{-1}}$, the matrix
$M_R[D,D]$ is symmetric: $M_R[v,w] = f_{w^{-1}v} = f_{v^{-1}w}
= M_R[w,v]$.  The matrix is also positive semidefinite (all
eigenvalues non-negative), since every eigenvalue of the staircase
element $\Pi$ on any irreducible representation is positive
\cite{Clio-EP}.

We define $M_R^{(k)}[D,D]$ analogously, replacing $\Pi$ by the
partial product $\Pi_k$.

%% ====================================================================
\section{Block decomposition}
%% ====================================================================

\subsection{Partial products and their matrices}

For $k = 0, 1, \ldots, n-1$, define $M_R^{(k)}[D,D]$ to be the
right-multiplication matrix of the partial staircase product $\Pi_k$,
restricted to $D$.  We have:
\[
  M_R^{(0)}[D,D] = I_{|D|} \qquad (\Pi_0 = e),
\]
\[
  M_R^{(n-1)}[D,D] = M_R[D,D] \qquad (\Pi_{n-1} = \Pi).
\]

\subsection{The block ratio}

At each step $k$, appending block~$k$ to the staircase replaces
$\Pi_{k-1}$ by $\Pi_k = B_k \cdot \Pi_{k-1}$ where
$B_k = \prod_{i=k}^{1} (1+s_i)/2$ is the product for block~$k$.
The corresponding matrices satisfy
\[
  M_R^{(k)}[v,w] = M_R^{(k-1)}[v,w] \cdot (\text{block-$k$ factor}),
\]
in the sense that $\det(M_R^{(k)}) = \det(M_R^{(k-1)}) \cdot R_k$
for some scalar $R_k$.

\begin{definition}[Block ratio]
The \emph{$k$-th block ratio} is
\[
  R_k \;=\; \frac{\det(M_R^{(k)}[D,D])}{\det(M_R^{(k-1)}[D,D])},
  \qquad k = 1, \ldots, n-1.
\]
This is well defined when $\det(M_R^{(k-1)}) \ne 0$.
\end{definition}

The block-multiplicative structure is the claim that each $R_k$ is a
positive integer.

%% ====================================================================
\section{Main theorem (partial result)}
%% ====================================================================

\begin{conjecture}[Block Multiplicativity]\label{conj:block-mult}
For all $n \ge 2$ and $k = 1, \ldots, n-1$:
\begin{enumerate}
  \item $\det(M_R^{(k)}[D,D])$ is a positive integer.
  \item The block ratio $R_k$ is a positive integer.
  \item Consequently, $\det(M_R[D,D]) = \prod_{k=1}^{n-1} R_k > 0$.
\end{enumerate}
\end{conjecture}

\begin{remark}
Conjecture~\ref{conj:block-mult} would complete the $H$-invariant
theorem: $\det(M_R[D,D]) > 0$ implies $M_R[D,D]$ is positive definite
(combining with positive semidefiniteness), which implies $D$ is a
basis for the image of $\Pi$, which by the total rank formula implies
$\rk(\Pi|_{V_\lambda}) = \dim V_\lambda^H$ for all $\lambda$.
\end{remark}

The theorem has been verified computationally for $n \le 6$; see
Section~\ref{sec:evidence}.

%% ====================================================================
\section{Computational evidence}\label{sec:evidence}
%% ====================================================================

All computations use exact rational arithmetic via SageMath/Python with
Young's seminormal form.  The staircase coefficients $f_g$ are computed
symbolically and the determinants evaluated as exact integers.

\subsection{Case $n = 4$}

$|D_4| = 6$, staircase word: $s_1,\; s_2 s_1,\; s_3 s_2 s_1$.
Three blocks.

\begin{center}
\begin{tabular}{cccc}
\toprule
Block $k$ & Partial product $\Pi_k$ & $\det(M_R^{(k)})$ & $R_k$ \\
\midrule
0 & $e$ & $1$ & --- \\
1 & $(1+s_1)/2$ & $1$ & $1$ \\
2 & $\Pi_1 \cdot (1+s_2)(1+s_1)/4$ & $24$ & $24$ \\
3 & $\Pi_2 \cdot (1+s_3)(1+s_2)(1+s_1)/8$ & $576$ & $24$ \\
\bottomrule
\end{tabular}
\end{center}

Thus $\det(M_R[D_4, D_4]) = 1 \cdot 24 \cdot 24 = 576 = 2^6 \cdot 3^2$.

\subsection{Case $n = 5$}

$|D_5| = 30$, four blocks.

\begin{center}
\begin{tabular}{cccc}
\toprule
Block $k$ & $\det(M_R^{(k)})$ & $R_k$ & Factored $R_k$ \\
\midrule
0 & $1$ & --- & \\
1 & $1$ & $1$ & $1$ \\
2 & $24^5$ & $24^5$ & $2^{15} \cdot 3^5$ \\
3 & $24^{10}$ & $24^5$ & $2^{15} \cdot 3^5$ \\
4 & $2^{70} \cdot 3^{10}$ & $2^{40}$ & $2^{40}$ \\
\bottomrule
\end{tabular}
\end{center}

Thus $\det(M_R[D_5, D_5]) = 24^{10} \cdot 2^{40} = 2^{70} \cdot 3^{10}$.

\medskip
\noindent
\textbf{Observation.}  The ratios $R_2 = R_3 = 24^5$ match the $n=4$
ratios $R_2 = R_3 = 24$, each raised to the power $|D_5|/|D_4| = 5$.
This is the tensor product phenomenon (Section~\ref{sec:tensor}).

\subsection{Case $n = 6$}

$|D_6| = 90$, five blocks.

\begin{center}
\begin{tabular}{cccc}
\toprule
Block $k$ & $R_k$ & Factored $R_k$ \\
\midrule
1 & $1$ & $1$ \\
2 & $24^{15}$ & $2^{45} \cdot 3^{15}$ \\
3 & $24^{15}$ & $2^{45} \cdot 3^{15}$ \\
4 & $r_H$ & $2^{91} \cdot 3^{11} \cdot 5^4 \cdot 7 \cdot 11^4$ \\
5 & $r_H$ & $2^{91} \cdot 3^{11} \cdot 5^4 \cdot 7 \cdot 11^4$ \\
\bottomrule
\end{tabular}
\end{center}

where we write $r_H = 2^{91} \cdot 3^{11} \cdot 5^4 \cdot 7 \cdot 11^4$
for the common value of $R_4$ and $R_5$.

The full determinant is:
\begin{align*}
  \det(M_R[D_6, D_6])
  &= 1 \cdot (24^{15})^2 \cdot r_H^2 \\
  &= 2^{272} \cdot 3^{52} \cdot 5^8 \cdot 7^2 \cdot 11^8,
\end{align*}
in agreement with the $q$-determinant computation \cite{Clio-qDet}.

\medskip
\noindent
\textbf{Symmetry $R_k = R_{n-k}$.}  In all cases computed, the block
ratios come in symmetric pairs: $R_1 = R_{n-1} = 1$ (boundary blocks),
and $R_k = R_{n-k}$ for all $k$.  This symmetry reflects the palindromic
structure of the staircase word under word reversal.

%% ====================================================================
\section{Tensor product structure}\label{sec:tensor}
%% ====================================================================

\subsection{The embedding $\Sym_4 \hookrightarrow \Sym_n$}

The sub-staircase involving only the generators $s_1, s_2, s_3$
(blocks 1, 2, 3 of the full staircase) is the complete staircase for
$w_0 \in \Sym_4$.  When $\Sym_4$ acts on the first four positions
inside $\Sym_n$, its staircase product is the restriction of $\Pi_3$
to $\Sym_4 \subset \Sym_n$.

\begin{proposition}[Tensor product structure, blocks 2 and 3]\label{prop:tensor}
Let $n \ge 4$ and let $\Pi_3^{(n)}$ denote the staircase partial
product after three blocks in $\Sym_n$.  Then
\[
  \det\!\bigl(M_R^{(3)}[D_n, D_n]\bigr)
  \;=\; \det\!\bigl(M_R^{(3)}[D_4, D_4]\bigr)^{|D_n|/|D_4|}.
\]
In particular, $R_2^{(n)} = (R_2^{(4)})^{|D_n|/|D_4|} = 24^{|D_n|/6}$
and $R_3^{(n)} = (R_3^{(4)})^{|D_n|/|D_4|} = 24^{|D_n|/6}$.
\end{proposition}

\begin{proof}[Proof (partial)]
The proposition is verified computationally for $n = 4, 5, 6$.
The underlying structure is:

The partial product $\Pi_3^{(n)}$ in $\Sym_n$ acts on $\QQ[\Sym_n]$
via right multiplication.  The right $\Sym_4$-coset decomposition
$\Sym_n = \bigsqcup_c c \cdot \Sym_4$ induces a corresponding
decomposition of $D_n$ into $|D_n|/|D_4|$ cosets of~$D_4$.

Since $\Pi_3^{(n)}$ involves only $s_1, s_2, s_3 \in \Sym_4$, its
right action decomposes across these cosets independently.
The matrix $M_R^{(3)}[D_n, D_n]$ is block-diagonal with
$|D_n|/|D_4|$ diagonal blocks, each isomorphic to
$M_R^{(3)}[D_4, D_4]$.  Hence its determinant is the
$|D_n|/|D_4|$-th power of $\det(M_R^{(3)}[D_4, D_4])$.

\medskip
\noindent
\emph{Note.}  The block-diagonal decomposition is not immediate from
the definitions and requires careful analysis of how the coset
structure interacts with the descent-paired condition.
A complete proof is not available.
\end{proof}

\begin{corollary}
The eigenvalues of $M_R^{(3)}[D_n, D_n]$ are exactly the eigenvalues
of $M_R^{(3)}[D_4, D_4]$, each with multiplicity $|D_n|/|D_4|$.
\end{corollary}

\noindent
Concretely: for $n = 5$, the matrix $M_R^{(3)}[D_5, D_5]$ of size
$30 \times 30$ acts as 5 independent copies of the $6 \times 6$ matrix
$M_R^{(3)}[D_4, D_4]$.  For $n = 6$, it acts as 15 copies.


%% ====================================================================
\section{Bridge pattern within blocks}
%% ====================================================================

\label{sec:bridge}

Within each block $k = (s_k, s_{k-1}, \ldots, s_1)$, the staircase
accumulates rank one generator at a time.  Define the \emph{partial
block-$k$ product}
\[
  B_k^{(j)} \;=\; \frac{(1+s_k)(1+s_{k-1})\cdots(1+s_j)}{2^{k-j+1}},
  \qquad j = k, k-1, \ldots, 1,
\]
so that $B_k^{(k)} = (1+s_k)/2$ is the first step and
$B_k^{(1)} = B_k$ is the full block.

Computational analysis for $n = 5, 6$ reveals the following pattern:

\begin{itemize}

\item \textbf{Even-indexed steps (generators $s_{2j}$):}  Appending a
factor $(1+s_{2j})/2$ at position $s_{2j}$ within the block causes a
rank drop of $|D|/6$ in the running restricted matrix.  (Here
$|D|/6$ is the ratio of the full descent-set to the $\Sym_4$ descent
set; for even $n$, this is an integer.)

\item \textbf{Odd-indexed steps (generators $s_{2j-1}$):}  Appending
a factor $(1+s_{2j-1})/2$ maintains or recovers rank.  No rank drop
occurs at these steps.

\item \textbf{Final step $s_1$:}  The final generator in every block
is $s_1 \in H$.  Appending $(1+s_1)/2$ always \emph{restores full
rank} within the block, undoing any transient drops.

\end{itemize}

\begin{remark}
This alternating rank pattern is consistent with the Rank Isolation
Lemma \cite{Clio-RankIso}, which states that rank drops in the full
staircase product occur \emph{only} at first appearances of
$H$-generators $s_1, s_3, s_5, \ldots$.  The bridge pattern describes
the finer within-block structure: temporary rank variations that cancel
by the end of each block.
\end{remark}


%% ====================================================================
\section{The perfect square structure}
%% ====================================================================

\begin{proposition}[Perfect square]\label{prop:square}
$\det(M_R[D,D])$ is a perfect square.  Specifically,
$M_R[D,D] = Q[D,:]^{\vphantom{T}} Q[D,:]^T$ where $Q$ is the
restriction to $D$ of a matrix $Q^{(n)}$ with $\Pi = Q Q^T$
in the orthogonal decomposition $L_\Pi = QQ^T$ on $\QQ[\Sym_n]$.
\end{proposition}

\begin{proof}[Proof (partial)]
The staircase product satisfies $\Pi = \Pi^2$ (modulo scaling) and
is a product of commuting projections in appropriate subgroups.
The factorization $L_\Pi = QQ^T$ arises from the eigenvalue
decomposition: since all eigenvalues of $M_R[D,D]$ are non-negative
(established in \cite{Clio-EP}), the Cholesky factorization exists
over $\QQ$, giving $M_R[D,D] = QQ^T$ with $Q$ having rational entries.
The perfect square property
\[
  \det(M_R[D,D]) \;=\; \det(Q[D,:])^2
\]
then follows from $\det(QQ^T) = \det(Q)^2$.

\noindent
\emph{Note.}  The claim that $Q$ can be chosen with integer entries
(explaining why $\det$ is a perfect square of an integer, not just
of a rational) is supported computationally but not fully proved.
\end{proof}

For $n = 4$: $\det = 576 = 24^2$, and indeed $\sqrt{576} = 24$.
For $n = 5$: $\det = 2^{70} \cdot 3^{10} = (2^{35} \cdot 3^5)^2$.
For $n = 6$: $\det = 2^{272} \cdot 3^{52} \cdot 5^8 \cdot 7^2 \cdot 11^8
= (2^{136} \cdot 3^{26} \cdot 5^4 \cdot 7 \cdot 11^4)^2$.

All verified perfect squares.


%% ====================================================================
\section{Open problem: proving $R_k > 0$}
%% ====================================================================

\begin{conjecture}[Positive block ratios]\label{conj:rk-pos}
For all $n \ge 2$ and $k = 1, \ldots, n-1$, the block ratio
$R_k = \det(M_R^{(k)}[D,D]) / \det(M_R^{(k-1)}[D,D])$ is a
positive integer.
\end{conjecture}

Proving Conjecture~\ref{conj:rk-pos} would yield:

\begin{corollary}[H-invariant theorem, conditional]
\label{cor:conditional}
Conjecture~\ref{conj:rk-pos} implies $\det(M_R[D,D]) > 0$, which
combined with positive semidefiniteness implies positive definiteness
of $M_R[D,D]$, which implies the $H$-invariant theorem
$\rk(\Pi|_{V_\lambda}) = \dim V_\lambda^H$ for all $\lambda \vdash n$
and all $n$.
\end{corollary}

\subsection{Possible approaches}

\paragraph{Approach 1: Induction on $k$ using the tensor structure.}
The tensor product structure (Proposition~\ref{prop:tensor}) handles
$R_2$ and $R_3$ for all $n$ at once, reducing them to the $n=4$ case.
A similar parabolic reduction might handle $R_4, R_5$ by reducing to
$n = 6$.  If each new pair $(R_{2j}, R_{2j+1})$ only depends on
$\Sym_{2j+2}$, the full result would follow by induction on $n$.

\paragraph{Approach 2: Hecke algebra $q$-deformation.}
The $q$-determinant $\Delta_n(q) = \det(M_R^q[D,D])$ is a polynomial
with all roots at $\mathrm{Re}(q) \le 0$ (verified $n \le 6$).
Combined with eigenvalue positivity \cite{Clio-EP}, this gives a
$q$-interpolation: if $\Delta_n(q) \ne 0$ for all $q > 0$, then
$\det(M_R^q[D,D]) > 0$ for all $q > 0$.  The block structure of
the $q$-determinant might factor as $\prod_k R_k(q)$ with each
$R_k(q) > 0$ for $q > 0$, giving $R_k = R_k(1) > 0$.

\paragraph{Approach 3: Kazhdan--Lusztig combinatorics.}
The matrix $M_R[D,D]$ has entries $f_{w^{-1}v}$ which are counts of
subwords of the staircase word equal to $w^{-1}v$.  The positivity
and integrality of $R_k$ might follow from a combinatorial
interpretation of the determinant in terms of non-intersecting lattice
paths or the permanent of a related matrix.


%% ====================================================================
\section{Failed approach: total non-negativity}
%% ====================================================================

\label{sec:tnn-fail}

The most natural approach to proving $\det(M_R[D,D]) > 0$ would be to
show that $M_R[D,D]$ is \emph{totally non-negative} (TNN): all minors
of all sizes are non-negative.  For TNN matrices, the determinant is
non-negative by definition.

\begin{proposition}[TNN fails]\label{prop:tnn-fail}
The matrix $M_R[D,D]$ is \emph{not} totally non-negative for $n \ge 5$.
\end{proposition}

\begin{proof}
Direct computation for $n = 5$: the $2 \times 2$ minor
$\det\!\begin{pmatrix} M_R[v_1, w_1] & M_R[v_1, w_2] \\
  M_R[v_2, w_1] & M_R[v_2, w_2] \end{pmatrix}$
is negative for specific choices of $v_1, v_2, w_1, w_2 \in D_5$.
\end{proof}

\begin{remark}
The TNN failure shows that the positivity of the full determinant is a
\emph{cancellation phenomenon}, not a consequence of uniform positivity
of all minors.  This makes a direct combinatorial proof more subtle:
the cancellations must be tracked carefully.

The failure already occurs at $n = 5$ with $2 \times 2$ minors,
indicating that TNN is not the right framework.  The correct
framework appears to be the eigenvalue positivity result
\cite{Clio-EP}, which shows that the negative minors cancel when
taking the full determinant.
\end{remark}


%% ====================================================================
\section{Summary of results}
%% ====================================================================

\begin{center}
\begin{tabular}{lll}
\toprule
\textbf{Statement} & \textbf{Status} & \textbf{Range} \\
\midrule
$M_R[D,D]$ is symmetric & Proved & All $n$ \\
$M_R[D,D]$ is PSD & Proved \cite{Clio-EP} & All $n$ \\
$\det(M_R[D,D]) > 0$ & Verified & $n \le 7$ \\
Block multiplicativity: $\det = \prod R_k$ & Verified & $n \le 6$ \\
$R_1 = R_{n-1} = 1$ & Verified & $n \le 6$ \\
$R_k = R_{n-k}$ (symmetry) & Verified & $n \le 6$ \\
$R_2 = R_3 = 24^{|D|/6}$ (tensor structure) & Verified + partial proof & $n \le 6$ \\
$\det$ is a perfect square & Verified & $n \le 6$ \\
$R_k > 0$ for all $k$ & \textbf{Open} & \\
TNN of $M_R[D,D]$ & \textbf{False} for $n \ge 5$ & \\
\bottomrule
\end{tabular}
\end{center}

\medskip\noindent
The central open problem is Conjecture~\ref{conj:rk-pos}.  If proved,
it would complete the $H$-invariant theorem for all $n$ by
Corollary~\ref{cor:conditional}.


\begin{thebibliography}{99}

\bibitem[Clio-EP]{Clio-EP}
Clio, \emph{Eigenvalue positivity for the Hecke staircase element
via Kazhdan--Lusztig positivity}, April 2026.

\bibitem[Clio-qDet]{Clio-qDet}
Clio, \emph{The $q$-determinant of the right multiplication submatrix
and the image basis conjecture}, April 2026.

\bibitem[Clio-RankIso]{Clio-RankIso}
Clio, \emph{Rank isolation and the staircase product: cascade
transport and parabolic reduction}, April 2026.

\bibitem[Clio-TotalRank]{Clio-TotalRank}
Clio, \emph{The total rank formula and image basis conjecture
for the staircase product}, April 2026.

\end{thebibliography}

\end{document}
